\scanpage{26}
% L1–6
\unsure{au|en|en un} sens technique suivant : Tout $X$-hom $G_{a,X}\mathrel{\longrightarrow}Y$ (où $Y$ est un préschéma sur $X$) qui est un $G_{a,X}$ hom (qud $G_{aX}$ opère sur le premier facteur grâce à $\Phi$, et sur $Y$ trivialement) est une \uline{section} de $Y$, i.e. se factorise en
\diagram{\begin{tikzcd}[ampersand replacement=\&]G_{aX}\arrow[r,"\mathrm{proj}"'] \& X\arrow[r] \& Y.\end{tikzcd}}
% L7–10
Cela se voit p.ex. sur le \deleted{fait que} \inserted{l'hom} d'\unsure{hyperalg.|hyperalgèbres} $\Phi$, en \unsure{notant|constatant|montrant} que celui-ci est injectif \margin{?}[du moins si $Y$ est affine\,; \deleted{\unsure{on peut}} \inserted{on peut} se ramener \unsure{tjrs} à ce cas].

% L11–12
% slanted
\margin{à justifier \unsure{théoriquement} ?}
On en conclut que l'on a \deleted{\unsure{que l'on a}} \inserted{une} suite exacte de cohomologie, qui s'écrit
% L13, cancelled beginning
\diagram{\begin{tikzcd}[ampersand replacement=\&,column sep=small]
\deleted{0}\arrow[r]\arrow[rr,dash,start anchor=west,end anchor=east,shift left=.35ex] \& \deleted{A[T]/T^p}\arrow[r]\arrow[loop above,draw=none,"{\rightarrow}"{description}] \& {}
\end{tikzcd}}
% L14
\diagram{\begin{tikzcd}[ampersand replacement=\&,column sep=4pt,font=\scriptsize]
0\arrow[r] \& H^0(X,\mathcal B_X)\arrow[r] \& A\arrow[r,"{\deleted{\unsure{x}}F}"] \& A\arrow[r] \& \unsure{H^1|H\!H}(X,B_X)\arrow[r] \& H^1(\unsure{Y|X},\underline O_X)\arrow[r,"F"] \& H^1(\unsure{X|Y},\underline O_X)
\end{tikzcd}}
% L15–18
[ ici $\Phi$ provient de Frobénius \deleted{\unsure{F}} \unsure{, et on|, est} \unsure{met|a mis} pour $A\colon H^0(X,\underline O_X)$ ] d'où $H^1(X,B_X)$ comme extension d'un \deleted{$A/A^p$} sous-groupe \deleted{\unsure{sous-groupe de}} \inserted{de $H^1(X,\underline O_X)$} par $A/A^p$. Si p.ex.
% L19–20
$X$ est une variété complète sur un corps algébriquement clos, alors on obtient pour
% archive sheet number: 25
